\documentclass[11pt,a4paper]{article}
\usepackage[T1]{fontenc}
\usepackage[utf8]{inputenc}
\usepackage{lmodern,microtype}
\usepackage[margin=25mm]{geometry}
\usepackage{amsmath,amssymb,amsthm,mathtools}
\usepackage{booktabs,longtable,array}
\usepackage{xcolor}
\usepackage{enumitem}
\usepackage{hyperref}
\hypersetup{colorlinks=true,linkcolor=blue!45!black,urlcolor=blue!55!black,
  citecolor=blue!45!black,pdftitle={Separated degrees and Pell quotients: smaller consistency equations for ZFC},
  pdfauthor={Research contribution prepared for ProveIt}}
\urlstyle{same}
\newcommand{\N}{\mathbb N}
\newcommand{\Z}{\mathbb Z}
\newcommand{\Pos}{\mathbb Z_{>0}}
\newcommand{\Con}{\operatorname{Con}(\mathrm{ZFC})}
\newcommand{\pop}{\operatorname{pop}}
\DeclareUrlCommand\code{\urlstyle{tt}}
\newtheorem{theorem}{Theorem}[section]
\newtheorem{lemma}[theorem]{Lemma}
\newtheorem{proposition}[theorem]{Proposition}
\newtheorem{corollary}[theorem]{Corollary}
\theoremstyle{definition}
\newtheorem{definition}[theorem]{Definition}
\setlength{\emergencystretch}{2em}
\setlength{\parskip}{3pt}
\setlist{itemsep=2pt,topsep=4pt}
\allowdisplaybreaks[2]

\title{Separated degrees and Pell quotients\\[4pt]
\large Smaller explicit consistency equations for ZFC}
\author{Research contribution prepared for Vladimir Reshetnikov's ProveIt project}
\date{23 October 2026}

\begin{document}
\maketitle
\begin{abstract}
We give several new explicit integer polynomials whose absence of zeros in
their stated domains is equivalent, by ordinary mathematics formalizable in
ZFC, to syntactic consistency of ZFC.  The strongest small-degree result has
twelve positive unknowns, degree \(242\), and a 94-operation division-free
evaluation circuit; an 89-operation alternative has degree \(258\).
Separating the ordinary degrees in a packed quadratic source yields a
nine-positive-unknown equation of degree \(2\,628\,944\) and 499 operations,
and a ten-unrestricted-integer equation of degree \(8\,317\,400\) and 544
operations.  A separate guard construction reduces a fifty-integer
quartic to forty-five integer unknowns.

The new ingredients are degree-separated coefficient extraction with an
unconditional positive carrier, two separate quadratic norms with exact
index recovery, a congruence-based quotient of a Pell ordinate, and shared
sign guards.  The native computation compiler, machine-to-refutation
calibration, and Jones/Sun arithmetic bridges are inherited and explicitly
identified.  We supply the polynomial circuits, fixed coefficient bytes,
mathematical proofs, independent finite audits, and source pins.  We claim
improvements over the inspected repository frontier, not an absolute
minimum, a decision of consistency, or a new end-to-end proof-assistant
certificate.
\end{abstract}

\section{Results and the meaning of simplicity}
\label{sec:results}

There is no single ordering of Diophantine equations that simultaneously
minimizes the number of unknowns, degree, arithmetic circuit size, coefficient
description length, and expanded monomial count.  The domain is also part
of the specification: positive unknowns cannot silently be treated as
unrestricted integers.  We therefore give explicit improvements in several
of these directions, and identify the precise comparisons.

The baseline is the repository revision
\begin{center}
\code{4b4349064a750ba6347293b55874a7f92abdd0db}.
\end{center}
It includes all eight inspected wave-7 archives under \code{docs/incoming}.
The comparison uses those archives as well as the integrated catalogue,
rather than comparing only against the older wave-6 summary.

\begin{table}[ht]
\centering\small
\begin{tabular}{@{}lrrl@{}}
\toprule
Domain and unknown count & Exact degree & Operations & Role\\
\midrule
\(\Pos^{12}\) & \textbf{242} & \textbf{94} & Lowest-degree new twelve-variable form\\
\(\Pos^{12}\) & \textbf{258} & \textbf{89} & Shorter twelve-variable evaluator\\
\(\Pos^{13}\) & 132 & 83 & Additional tradeoff\\
\(\Pos^9\) & \textbf{2,628,944} & \textbf{499} & Nine positive unknowns\\
\(\N^9\) & 2,628,944 & 501 & Explicit nonnegative translation\\
\(\Z^{10}\) & \textbf{8,317,400} & \textbf{544} & Ten unrestricted integers\\
\(\Z^{45}\) & \textbf{4} & --- & Quartic, 28 quadratic residuals\\
\bottomrule
\end{tabular}
\caption{Delivered constructions.  Here \(\N=\{0,1,2,\ldots\}\).
Operations are live binary additions, subtractions, and multiplications.}
\label{tab:results}
\end{table}

The incoming twelve-positive form has degree 897 and 108 operations.
Degree 242 is a \(73.02\%\) reduction; the 89-operation variant reduces
its evaluator by \(17.59\%\).  The best incoming nine-positive degree is
\(7\,344\,328\), while the smallest incoming operation count at nine
positive variables is 548, attained by a different equation.  Our equation
improves both independent minima: the degree falls by \(64.20\%\) and
the operation count by \(8.94\%\).  Similarly, the incoming ten-integer
minima are degree \(21\,111\,400\) and 582 operations, again at different
points.  The new degree falls by \(60.60\%\) and the operation count by
\(6.53\%\).  The quartic result decreases the unrestricted variable count
from 50 to 45, a \(10\%\) reduction.  The detailed archive-by-archive
comparison is retained in \code{proofs/WAVE7_AUDIT.md}.

These changes do not improve every existing objective.  In particular,
the known 31-positive-variable quartic and short fifteen-positive-variable
evaluators remain useful points on the frontier.  We do not reduce the
number nine of positive unknowns or ten of unrestricted unknowns.  The
forty-five-variable quartic has more expanded symbolic terms than its
chosen parent, as discussed in Section~\ref{sec:quartic45}.

\begin{theorem}[Mathematical endpoint]
For every polynomial in Table~\ref{tab:results}, with its domain \(D\)
and fixed coefficient instance specified in this packet, the supplied
reductions and the identified inherited compiler/calibration theorems give
an ordinary mathematical proof formalizable in ZFC of
\[
  \exists\mathbf x\in D\;P(\mathbf x)=0
  \quad\Longleftrightarrow\quad
  \text{there exists a finite syntactic ZFC refutation}.
\]
Consequently ZFC proves
\[
 \boxed{\quad
 (\forall\mathbf x\in D\;P(\mathbf x)\ne0)
 \quad\Longleftrightarrow\quad\Con.
 \quad}
\]
\end{theorem}

The theorem is a composition of explicit new reductions with named inherited
theorems.  Its mathematical dependency is the accepted K215 consistency
calibration, not an assertion that bounded simulations establish nonhalting.
The formalization status is narrower: no newly assembled derivation of
the displayed equivalence in the repository's object-language \code{BProv}
calculus is supplied.  Section~\ref{sec:status} explains this distinction.

For a practical choice, the degree-242 twelve-positive form is the most
balanced new construction in this packet.  If unknowns must range over all
integers, use the ten-integer form for a small variable count or the
forty-five-integer form for degree four.  Every register appearing below
is a polynomial abbreviation; only the explicitly listed coordinates
are quantified.

\section{Explicit coefficients and inherited interfaces}
\label{sec:coefficients}

The nine-, ten-, twelve-, and thirteen-variable constructions use one
complete coherent four-port tuple
\[
  (\beta,K_0,M_C,M_F),\qquad \beta=2^{590069}-1.
\]
Its identifier is
\begin{center}\small
\code{wang943-checksum67-shared-anchor-length8807-prime581197}.
\end{center}
The bit lengths of its four ports are respectively
\((590069,592884,590069,590070)\).  All four unsigned little-endian
integers are included under \code{coefficients/coefficients/}, together
with their construction and an independent decoder.  They occupy 295,388
raw bytes and 9,445 gzip bytes.  The fourth port \(M_F\) is already the
shifted \(M_{F,\mathrm{native}}+2^{590069}-1\); it must not be shifted
a second time.  The archive contains the literal 943-tile input and the
literal K215 transition input needed by its builder.

The raw SHA-256 hashes are given in \code{COEFFICIENT_BINDING.json}.
The independent decoder reconstructs all four whole integers, verifies
66,062 native basis identities, and checks the actual coefficient residues
needed by the compiler's rank controller.  Its 38,663 one-bit negative
controls supplement these exact checks.  The unbounded compiler proof is
included separately; finite checks are not a replacement for it.

The forty-five-integer quartic uses the complete package-34 tuple
\begin{center}\code{mixed-parity-unrounded-1064921}.\end{center}
It has a different binding, even
though the arithmetic register names are similar.  Its four raw payloads
occupy 532,479 bytes.  Keeping that complete tuple is necessary for its
odd-exponent canonical radix.  No coefficient ports are spliced between
the two instances.

For the degree-separated source below, all channel coefficient recipes
are evaluated at the first tuple.  A termwise absolute bound gives
\[
 M=2^{1182955},\qquad \kappa=1182968,
 \qquad a_S=2^\kappa=8192M.
\]
Every channel coefficient has absolute value strictly below \(M\).
The two additional scale coefficients are
\(2^{9463734}\) and \(2^{1182947}\).  These scales and every derived
channel coefficient leaf are also materialized in \code{fixed/}, with
a signed magnitude format and exact hashes.  There are fourteen such
nonliteral leaves, requiring 2,512,713 raw bytes or 47,720 gzip bytes.
This byte count is separate from arithmetic circuit complexity.

The inherited interfaces are the positive quadratic native compiler,
the normalized Pell compiler and its signed auxiliary-index lemma,
the accepted K215-to-refutation calibration, the Jones nine-variable
ratio bridge, and the Sun ten-integer bridge.  The relevant mathematical
proof texts and Jones Lean source are retained under \code{upstream/},
\code{low_degree/upstream/}, and \code{quartic/upstream/}.  Their source
revision and byte hashes are recorded.  These imports are part of the
proof, not hidden assumptions about the existence of a solution.

\section{Two separate norms and a balanced Pell quotient}
\label{sec:low-degree}

The direct positive-variable route admits a substantial reduction that is
independent of the nine-variable packing construction.  The recommended
polynomial has twelve positive unknowns, exact degree $242$, and $94$ binary
operations.  A cheaper companion has degree $258$ and $89$ operations.
The incoming wave-7 comparison is twelve positive unknowns, degree $897$,
and $108$ operations.  A further thirteen-positive companion has degree
$132$ and $83$ operations.  The twelve-positive degree-$314$, $91$-operation
intermediate is retained because its reconstruction makes the mechanism
especially transparent.

The new mechanism keeps two quadratic norms separate.  One combines the
main Pell root and the first Pell root while absorbing the index quotient.
The other removes the strong Pell root.  A known residue of the auxiliary
Pell ordinate then permits an exact polynomial factor to be removed, with
its size chosen to balance the two remaining terms.  All coefficients are
the complete compact tuple in the coefficient appendix, with the already
shifted high mask.  There are no additional existential coefficient ports.

\subsection{The literal twelve-variable polynomial}

Supply exactly
\[
 J,F,\alpha,z,i,T,s,w,\eta,\zeta,y,t\in\mathbb Z_{>0}.
\]
Every other letter in this section denotes a polynomial abbreviation.  Set
\begin{align*}
 q&=\beta J+1,&X&=wq,&Y&=sq^3,&E&=XY,\\
 k&=\eta+\zeta,&c&=kY+\eta,&a&=E+Y,&H&=4a+3,\\
 \Delta&=a^2+H,&\ell&=EkY,&A&=1+\ell(\ell+k),&G&=1+\Delta c^2,\\
 b_0&=X+ac,&U&=4q^4,&r&=k-R-1,&S&=UH.
\end{align*}
Here the transport expressions are
\begin{align*}
 R&=((q-1)(q-F)+\alpha)(q^2-1)+(M_C+qM_F)J,\\
 N_t&=(K_0+w)(q-F-\alpha)+q-F-z(q-1).
\end{align*}
Define the first norm by
\begin{align}
 W&=Ub_0+H\bigl(U(t+\ell)+r\bigr),\notag\\
 \mathcal N&=(W^2+U^2G-S^2A)^2-4W^2U^2G.
 \label{eq:two-root-norm}
\end{align}
For the second norm put
\begin{align*}
 I&=ic^2,&D&=1+\Delta I^2,&Q&=(\Delta I)^2,&d&=Q-1,\\
 v&=cT,&b&=c+RD.
\end{align*}
Choose one of the following two pairs of polynomial abbreviations:
\[
 (M,C)=\bigl((\Delta ic)^2,c^2\bigr)
 \quad\hbox{or}\quad
 (M,C)=(Q,1).
\]
In both cases $Q=MC$.  Finally set
\begin{align}
 Z&=C(Dv^2+b^2-1)+dy(2-My),\notag\\
 \mathcal J_M&=Z^2-4C^2Dv^2b^2,\notag\\
 \boxed{P_M=N_t(1+\mathcal N)(1+\mathcal J_M)-1=0.}
 \label{eq:twelve-balanced}
\end{align}
The first choice is the degree-$242$ equation; the second is the
$258$-degree companion.  Both are ordinary integer polynomials.  The
balanced evaluator computes $M=(\Delta I)(\Delta i)$ and $Q=Mc^2$,
which is the same polynomial definition with one multiplication saved.

\begin{theorem}
For each of the two fixed polynomials in \eqref{eq:twelve-balanced}, ZFC
proves that it has no solution in twelve strictly positive integers if and
only if $\operatorname{Con}(\mathrm{ZFC})$.  Their exact degree/operation
pairs are $(242,94)$ and $(258,89)$, respectively.
\end{theorem}

The theorem composes the reduction below with the inherited native-window
compiler and K215 syntactic-refutation calibration.  The complete new
argument is supplied here; the source texts for the signed Pell recovery,
canonical compiler witnesses, and auxiliary congruences are retained in
\texttt{low\_degree/upstream/}, with exact pins.  This is an ordinary
mathematical proof formalizable in ZFC, rather than a new exported
object-language derivation or a newly checked end-to-end proof-assistant
artifact.

\subsection{Units, positivity, and recovery of the strong root}

Both $\mathcal N$ and $\mathcal J_M$ are squares modulo $4$.  If
\eqref{eq:twelve-balanced} holds, its three integer factors are units.
The value $-1$ for either $1+\mathcal N$ or $1+\mathcal J_M$ would force
the corresponding norm to equal $-2$, impossible modulo $4$.  Thus
\begin{equation}
 N_t=1,\qquad \mathcal N=\mathcal J_M=0.
 \label{eq:low-units}
\end{equation}
This obtains the transport interface before any Pell index argument.  Its
inherited bounds are
\begin{equation}
\begin{gathered}
 q\ge16,\quad C_w=q-F-\alpha>0,\quad 3q<R<q^4-2,\\
 X\ge q,\quad Y\ge q^3,\quad E\ge q^4>R+2,\qquad
 k\ge2,\quad kY<c<k(Y+1).
\end{gathered}
\label{eq:low-outer}
\end{equation}
For example, $C_w\le0$ would imply
$N_t\le q-F-z(q-1)=1-F-(z-1)(q-1)\le0$.  The remaining bounds are the
authentic whole-tuple mask bounds.  Consequently $M,C,v,b$ are positive,
$Q>1$, and $M\ge2$.  All denominators used next are nonzero.

From $\mathcal J_M=0$ define the rational number
\[
 f=\frac{Z}{2Cvb}.
\]
It satisfies $f^2=D$, so it is an integer.  The sign of this recovered
integer is not prescribed.  Put $y_{\rm old}=My-1>0$.  The exact polynomial
identity
\begin{equation}
 MZ=Q(Dv^2+b^2)-1-d(My-1)^2
 \label{eq:quotient-identity}
\end{equation}
then gives
\begin{equation}
 f^2-\Delta I^2=1,\qquad
 d y_{\rm old}^2=Q(vf-b)^2-1,
 \qquad vf-b=c(Tf-1)-Rf^2.
 \label{eq:low-strong-aux}
\end{equation}
These are the normalized strong and auxiliary Pell equations.  The signed
auxiliary theorem uses $|f|$ as its modulus and covers a negative $f$.
To match the preserved ordinary-strong statement explicitly, take
$i_{\rm ordinary}=\Delta i>0$; then
$(i_{\rm ordinary}c^2)^2=\Delta(f^2-1)=Q$.  No positive sign for the
reconstructed $f$ is silently assumed.

\subsection{A two-root carrier restores the main square}

Write $g_0=\sqrt G>0$ and $\tau_0=\sqrt A>0$.  There is the exact real
factorization
\begin{equation}
 \mathcal N=\prod_{\sigma,\epsilon\in\{-1,1\}}
       (W-\sigma Ug_0-\epsilon S\tau_0).
 \label{eq:main-four-factors}
\end{equation}
Since $r=k-R-1\ge1-R$, $U>R$, and $t\ge1$, we have $t+r/U>0$ and
\[
 W=Ub_0+S(\ell+t+r/U)>0.
\]
A vanishing factor supplies real roots $g^2=G$, $\tau^2=A$ with
$W=Ug+S\tau$.  Squaring isolates
\[
 2WU g=W^2+U^2G-S^2A.
\]
Thus $g$ is rational, hence integral; then $\tau$ is integral as well.
This argument remains valid when the two radicands have the same square
class.

The sign $g<0$ is impossible.  Indeed
\[
 \tau_0<\ell+k,\qquad
 \frac{b_0+g_0}{H}>\frac{2ac}{4a+3}>\frac c3>k,
\]
since $a>1$ and $Y>3$.  For $g=-g_0$ the factor equation would give
\[
 \tau=\ell+t+r/U+(b_0+g_0)/H>\ell+k>\tau_0,
\]
a contradiction.  Therefore $g=g_0>0$.  Reducing the carrier identity
modulo $U$ proves
\begin{equation}
 U\mid H(k-R-1).
 \label{eq:low-weak-divisibility}
\end{equation}
We have not yet canceled $H$ or recovered the omitted index quotient.

\subsection{Recover the exact index before the deleted quotients}

Put $A_0=a+2$ and $P_0=2EY+1$.  Standard Pell classification gives
positive indices $p,m,n$ with
\begin{align*}
 g&=\chi_{A_0}(p),&c&=\psi_{A_0}(p),\\
 |f|&=\chi_{A_0}(m),&I&=\psi_{A_0}(m),\\
 |\tau|&=\chi_{P_0}(n),&k&=2\psi_{P_0}(n).
\end{align*}
The first-Pell discriminant is $EY(EY+1)$, with fundamental
integer-coordinate unit $(P_0,2)$.  Because $P_0>A_0$ and $c>k$,
monotonicity gives $n<p$.  Also $c>1$ implies $p\ge2$, whence
$c\ge2A_0>2R$, and Pell growth gives $c>2p$.

Strong divisibility and $c^2\mid I$ give $p\mid m$.  Write $m=pj$.
Then
\[
 I=c\psi_g(j),\qquad \psi_g(j)\equiv jg^{j-1}\pmod c.
\]
Since $g^2\equiv1\pmod c$, $g$ is invertible modulo $c$, so $c\mid j$.
Consequently $c\mid m$ and $m\ge pc>2p$.

For clarity, the signed auxiliary-index recovery used at this point does
not depend on either deleted quotient.  Put $M_0=\Delta I$ and
$V=c(Tf-1)-Rf^2$.  The second equation in \eqref{eq:low-strong-aux} is
\[
 M_0^2V^2-(M_0^2-1)y_{\rm old}^2=1.
\]
Its Pell index $j_0$ is odd because its abscissa is divisible by $M_0$.
For the integer odd-index polynomials $Q_j$,
\begin{align*}
 \chi_{M_0}(2j+1)&=M_0Q_j(M_0^2),\\
 Q_j(1-A_0^2)&=(-1)^j\psi_{A_0}(2j+1),\\
 Q_j(0)&=(-1)^j(2j+1).
\end{align*}
Modulo $|f|$, $M_0^2\equiv1-A_0^2$ and $V\equiv-c$.  Squaring and
doubling, and reducing the signed abscissa index using $0<2p<m$, yields
$j_0\equiv\pm p\pmod m$, hence modulo $c$.  Modulo $c$,
$M_0^2\equiv0$, $f^2\equiv1$, and $V\equiv-R$, so
$j_0\equiv\pm R\pmod c$.  Since $0<p,R<c/2$, it follows that
\begin{equation}
 p=R.
 \label{eq:low-index}
\end{equation}
This is the signed index argument reproduced in the preserved sources;
its inputs have all been verified before using $h$ or $\gamma$.

Let $\mu=U/\gcd(U,H)$.  Since $H$ is odd and $H\equiv3\pmod q$, the
only possible common prime of $q$ and $H$ is $3$.  If $3\mid q$, then
$27\mid Y$ and $H\equiv3\pmod{27}$, so the $3$-adic valuation is one.
Therefore
\[
 \gcd(U,H)\le3,\qquad \mu\ge4q^4/3>R.
\]
The first Pell recurrence gives $k\equiv2n\pmod{4EY}$.  The modulus
$4EY=4ws^2q^7$ is a multiple of $U$.  Thus
\eqref{eq:low-weak-divisibility} gives $2n\equiv R+1\pmod\mu$.
But $1\le n<R$ implies $|2n-R-1|<R<\mu$.  Hence
\[
 2n=R+1,\qquad r=k-R-1\equiv0\pmod{4EY}.
\]
In particular $h=r/E$ and $r/U$ are integers.  Moreover $n\ge25$ and
$\psi_{P_0}(n)>n$ make $h>0$.  Dividing the carrier by $U$ gives
\[
 g=b_0+H(t+\ell+r/U-\tau),
\]
so $\gamma=(g-b_0)/H$ is integral.  Since $n\ge2$ gives
$k\ge4P_0>X$, we have $c>X$; since $\Delta>(a+1)^2$,
$g-b_0>c-X>0$.  Thus $\gamma>0$.

All normalized compiler equations have now been recovered, including
$k-hE-R=1$ and the positive main quotient.  Their signed auxiliary
soundness theorem gives a blank K215 halt, hence a finite syntactic ZFC
refutation.  The proof does not require making the reconstructed $f$
positive on the same assignment.

\subsection{Canonical completeness and the removed Pell residue}

Start with a canonical positive normalized compiler witness.  It has
$p=R=2n-1$, $h>0$, and $\gamma>0$.  Its first Pell recurrence gives
$U\mid r$.  Define
\begin{equation}
 t=\gamma+\tau-\ell-r/U.
 \label{eq:low-t-complete}
\end{equation}
This is a positive integer: since $\tau^2=\ell^2+\ell k+1$ and
$\ell>k$,
\[
 \tau-\ell=\frac{\ell k+1}{\tau+\ell}
 >\frac{\ell k}{2\ell+k}>\frac k3>\frac rU.
\]
The carrier identity follows, so $\mathcal N=0$.

For the new ordinate quotient, one may replace the canonical auxiliary
witnesses explicitly.  The authentic population decoding gives
$R\equiv3\pmod4$.  To identify the precise inherited fact, the native
radix is divisible by $4$, $q=B^N$, $J\equiv1\pmod4$, and
$M_C\equiv2\pmod4$.  The low mask forces $C_w\equiv1\pmod4$; the
formula for $R$ gives $R\equiv C_w+2\equiv3\pmod4$.  This argument is
preserved in the wave-7 squared-modulus source and applies to the whole
compact tuple here.

Keep the outer coordinates and the original positive $f,i$.  Put
$m_0=\Delta I$, so $Q=m_0^2$, and choose
\begin{equation}
 V=\chi_{m_0}(R)/m_0,\qquad
 y_{\rm old}=\psi_{m_0}(R).
 \label{eq:canonical-auxiliary}
\end{equation}
The odd-index polynomial makes $V$ a positive integer.  Its identities give
\begin{equation}
 V\equiv-c\pmod{f^2},\qquad
 V\equiv-R\pmod c,
 \qquad y_{\rm old}\equiv-1\pmod Q.
 \label{eq:canonical-residues}
\end{equation}
For the first use $m_0^2\equiv1-A_0^2\pmod{f^2}$ and
$(R-1)/2$ odd; for the second use $c\mid m_0$ and the constant term
of $Q_{(R-1)/2}$.  For the third, $\psi_x(R)$ is an even polynomial
in $x$, with constant term $(-1)^{(R-1)/2}=-1$.  The recurrence
$\psi_x(j+1)=2x\psi_x(j)-\psi_x(j-1)$ proves this assertion by induction.

Because $\gcd(c,f)=1$ and $f^2\equiv1\pmod c$, the numerator below is
divisible by $cf$ (in fact by $cf^2$).  Thus
\[
 T=\frac{V+c+Rf^2}{cf}>0
\]
is an integer, and $V=c(Tf-1)-Rf^2$.  The Pell identity gives the auxiliary
equation.  This replaces only $T$ and the auxiliary ordinate; it leaves
\eqref{eq:low-t-complete} unchanged.  No assumption of cofinal auxiliary
witnesses is needed.

For either selected $M\mid Q$, define
\[
 y=(y_{\rm old}+1)/M\in\mathbb Z_{>0}.
\]
Identity \eqref{eq:quotient-identity} gives $\mathcal J_M=0$.  All twelve
supplied coordinates are positive, proving completeness.  Negating the
existence equivalence proves the stated consistency equivalence in ZFC.

\subsection{Exact degrees and the intermediate companions}

Write the leading homogeneous forms as
\begin{gather*}
 q_0=\beta J,\quad X_0=wq_0,\quad Y_0=sq_0^3,\quad E_0=X_0Y_0,\quad
 k_0=\eta+\zeta,\quad c_0=k_0Y_0,\quad \ell_0=E_0c_0,\\
 H_0=4E_0,\quad U_0=4q_0^4,\quad I_0=ic_0^2,\quad
 R_0=q_0^3(q_0-F),\quad (N_t)_0=w(q_0-F-\alpha)-zq_0.
\end{gather*}
The large apparent degree of \eqref{eq:two-root-norm} cancels by an exact
polynomial identity.  Set
\begin{align*}
 D_0&=b_0^2-G=X^2+2Xac-Hc^2-1,&Z_0&=Ut+r,\\
 C_2&=U^2D_0+2UHb_0Z_0
       +H^2(2U\ell Z_0+Z_0^2-U^2\ell k-U^2).
\end{align*}
Then
\begin{equation}
\begin{split}
 \mathcal N={}&4U^4H^2\ell^2D_0+4U^2Hb_0\ell C_2+C_2^2\\
 &\hspace{1cm}-4U^4H^2G(\ell k+1).
\end{split}
\label{eq:norm-degree-identity}
\end{equation}
Here $\deg D_0\le16$ and $\deg C_2\le32$.  The four summands have
respective degree bounds $66,68,64,62$.  Only the second reaches $68$,
with nonzero highest part
\begin{equation}
 \mathcal N_{68}=2^{16}w^6s^9k_0^3q_0^{49}(2t-k_0).
 \label{eq:low-norm-leader}
\end{equation}
This proves the upper and lower degree statements by polynomial identities,
without treating a sampled diagonal degree as an upper bound.

For the full-modulus choice $M=Q$, the two summands defining $Z$ have
maximum degrees $76$ and $94$.  The latter has highest part
$-Q_0^2y^2$, where $Q_0=E_0^4I_0^2$.  Thus $\mathcal J_M$ has degree
$188$, and $P_M$ has degree $68+188+2=258$, with highest part
\[
 \mathcal N_{68}Q_0^4y^4(N_t)_0\ne0.
\]
For the balanced choice, $\deg M=36$ and $\deg C=10$.  The first summand
of $Z$ has degree $86$ and highest part $c_0^2R_0^2E_0^4I_0^4$; the
second has degree $84$.  The subtracted term of $\mathcal J_M$ has degree
$142$.  Hence $\deg\mathcal J_M=172$, and the highest part of $P_M$ is
\[
 \mathcal N_{68}c_0^4R_0^4E_0^8I_0^8(N_t)_0\ne0,
\]
of degree $242$.

The degree-$314$ intermediate uses the original ordinate directly:
\begin{align*}
 e&=Q(Dv^2+b^2)-1,& B_{\mathrm{aux}}&=2Qvb,\\
 \mathcal J&=(e-dy^2)^2-B_{\mathrm{aux}}^2D,&
 P_{314}&=N_t(1+\mathcal N)(1+\mathcal J)-1.
\end{align*}
It reconstructs $f=(e-dy^2)/B_{\mathrm{aux}}$, with no canonical residue
restriction.  Its inner norm has degree $244$ and highest part
$E_0^{16}I_0^{12}R_0^4$, proving total degree $314$.

Finally, retaining $f>0$ gives the thirteen-positive companion
\begin{align*}
 N_s&=f^2-\Delta I^2,&V&=c(Tf-1)-Rf^2,\\
 N_a&=\Delta(f^2-1)(V^2-y^2)+y^2,&
 P_{132}&=N_t(1+\mathcal N)N_sN_a-1.
\end{align*}
A zero forces all factors to be integer units and $\mathcal N=0$.
The elementary negative-Pell obstruction
\[
 x^2-(m^2-1)y^2\ne-1\qquad(m\ge2)
\]
first rules out $N_s=-1$.  Then $\Delta(f^2-1)=(\Delta I)^2$ is a
square, so the same obstruction rules out $N_a=-1$.  Thus $N_t=N_s=N_a=1$.
For the obstruction, a hypothetical nonnegative solution has
$(m-1)y<x<my$; multiplying by $m-\sqrt{m^2-1}$ gives a smaller positive
ordinate $my-x$, contradicting descent.  The remaining proof is the
same carrier/index reconstruction above.  Degrees of the factors are
$2,68,34,28$, and their nonzero leading product is
\[
 -\mathcal N_{68}E_0^4I_0^2c_0^2T^2f^4(N_t)_0,
\]
so the exact degree is $132$.

The literal DAGs have $94$, $89$, $91$, and $83$ live binary gates for
$P_{242},P_{258},P_{314},P_{132}$, respectively.  Independent finite checks
include full polynomial arrays over three primes using the actual fixed
coefficients, four-factor multiplication in the formal quadratic algebra,
unit-guard congruences, the exact quotient identity, and nontrivial canonical
Pell reconstructions.  These validate the delivered arithmetic; the universal
proof and the identified inherited compiler arguments establish the theorem.

\section{One packed word, three ordinary-degree channels}
\label{sec:channels}

The incoming selective encodings improve on full Sidon separation by
isolating only monomials that actually occur in the parent system.  Their
homogenized packed square still mixes constant, linear, and quadratic
source terms.  The new observation is that these three ordinary degrees
can be evaluated separately without introducing another unknown.

Let \(r_1,\ldots,r_{22}\) be the raw quadratic residuals of the
31-positive-variable parent in \code{data/parent_quartic31.json}.
The raw positive coordinates are retained, rather than expanding a
positive-to-nonnegative shift into the source coefficients.  In their
specified order, use the weights
\[
\begin{split}
(w_1,\ldots,w_{31})={}&(104,76,79,60,108,8,99,55,36,54,98,11,7,10,44,32,\\
 &106,22,17,1,3,4,37,89,97,107,39,40,86,45,92),
\end{split}
\]
and, in residual order, the target positions
\[
\begin{split}
(q_1,\ldots,q_{22})={}&(1081,1388,1366,763,850,546,1414,1147,1914,1093,821,\\
 &1861,39,1621,238,759,500,1711,1052,283,1343,368).
\end{split}
\]
Put
\[
 W=108,\quad Q_c=1914,\quad K_c=Q_c+2W+1=2131,
 \quad E_c=K_c+W+2=2241,
\]
\[
 H_c=E_c+W+23=2372.
\]
The 31-coordinate order is exactly
\[
\begin{split}
&(J,F,\alpha,z,f,h,i,o,j,s,w,\tau,\eta,\zeta,y,\gamma,\\
&\hspace{13mm}N,E,P,c,a,X,D,C_2,L,F_2,V,g,I,B_{\rm aux},\Delta).
\end{split}
\]

For a source coefficient \(p_{j,\alpha}\), write \(d=|\alpha|\).
Place its reversed term at exponent \(q_j-\alpha\cdot w\) in
\(A_d(T)\).  Its coefficient is \(2p_{j,\alpha}\) for a constant,
a linear monomial, or a square; for a product of two different coordinates
it is \(p_{j,\alpha}\).  Appendix~\ref{app:channels} gives every
coefficient and exponent.  There are 45 quadratic records, 41 linear
records, and eight constant records.

\begin{lemma}[Degree-separated extraction]
For indeterminates \(z_1,\ldots,z_{31}\), define
\[
 Y_z(T)=\sum_i z_iT^{w_i},\qquad
 \mathcal F(T,y)=A_2(T)y^2+A_1(T)y+A_0(T)+T^{Q_c+1}(y^2+1).
\]
Then, for every \(1\le j\le22\),
\[
 [T^{q_j}]\mathcal F(T,Y_z(T))=2r_j(z_1,\ldots,z_{31}).
\]
\end{lemma}
\begin{proof}
The finite certificate checks, separately for each \(d\in\{0,1,2\}\),
that for every supported monomial \(\alpha\) of degree \(d\) and every
monomial \(\delta\) of exactly that degree,
\[
 q_j-\alpha\cdot w+\delta\cdot w=q_k
 \quad\Longleftrightarrow\quad (j,\alpha)=(k,\delta).
\]
All reversed exponents are nonnegative.  This is verified against all
496 pure quadratic monomials, all 31 linear monomials, and the constant
monomial: \(45\cdot496+41\cdot31+8=23599\) interactions.  The
coefficient of a mixed product in \(Y_z^2\) is two, while that of a
square is one.  The stated coefficient factors give exactly \(2r_j\).
All carrier exponents exceed \(Q_c\).  The supplied checker also
reconstructs the 22 identities symbolically with the native ports left
indeterminate, independently of the interaction count.
\end{proof}

Cross-degree collisions are harmless because a term of \(A_1Y_z\)
is never multiplied by a quadratic monomial, and a term of \(A_2Y_z^2\)
is never multiplied by a constant or a linear monomial.  The three
channels are fixed polynomials evaluated at the same base and the same
single packed input.  No extra existential coordinate is introduced.

\subsection{A positive carrier valid before decoding}

Let \(b\ge2\), \(B=a_Sb^2\), and
\[
 \ell=\sum_iB^{w_i},\qquad y=g+\ell.
\]
The shift by \(\ell\) means that packed nonnegative digits \(n_i\)
decode to the positive parent coordinates \(n_i+1\).  By construction,
the reversed exponents are distinct within each channel.  Since each
coefficient has magnitude below \(M\),
\[
 |A_d(B)|<\frac{MB^{Q_c+1}}{B-1}.
\]
Put \(\varepsilon=M/(B-1)<1/4\).  The elementary inequality
\(y^2+|y|+1\le\tfrac32(y^2+1)\) gives, for \emph{every integer}
\(y\),
\begin{equation}
 \tfrac12B^{Q_c+1}(y^2+1)
 <\mathcal F(B,y)<2B^{Q_c+1}(y^2+1).
 \label{eq:carrier-positive}
\end{equation}
In particular, positivity is available before a ratio bridge has recovered
a bound on \(g\).  A carrier proved positive only on the eventual
decoding range would not justify that proof order.

Define the source registers
\begin{align*}
 N_0&=B^{W+1},& Z_c&=B^{K_c+1},& A_T&=\sum_jB^{q_j},\\
 M_c&=N_0-1-(b-1)\ell,& J_c&=2\mathcal F(B,g+\ell)+A_T,\\
 S_c&=g+N_0J_c,&T_c&=M_c+N_0(B-2)A_T.
\end{align*}
For every \(g>0\), \(S_c>0\).  The base-\(B\) digits of \(M_c\)
are \(B-b\) at the 31 weights and \(B-1\) at the remaining positions
through \(W\), so \(0\le M_c<N_0\) and \(T_c\ge0\), without any
power-of-two assumption on \(b\).

If \(0\le g<4bB^W\), then \(g<N_0\), \(\ell<2B^W\), and
\(g+\ell<5bB^W\).  Equation~\eqref{eq:carrier-positive} gives
\[
 0<\mathcal F(B,g+\ell)<52b^2B^{K_c},\qquad
 0<J_c<105b^2B^{K_c}<Z_c.
\]
Also \((B-2)A_T<Z_c\).  Thus \(0\le S_c,T_c<B^{E_c}\).
The stronger bound needed for the ten-integer adapter follows from
\[
 J_c+(B-2)A_T
 =2\mathcal F(B,g+\ell)+(B-1)A_T
 <(104b^2+1)B^{K_c}<\tfrac14Z_c.
\]
Here \(a_S>420\).  Since \(g+M_c<2N_0\) and \(Z_c\ge4\),
\begin{equation}
 S_c+T_c<N_0(2+Z_c/4)<N_0Z_c=B^{E_c}.
 \label{eq:source-sum-range}
\end{equation}

\subsection{Exact zero extraction with signed borrows}

Assume \(b\) is a power of two.  Then \(B\) and \(N_0\) are powers
of two.  In the range just proved,
\[
 g\mathbin{\&}M_c=0
 \quad\Longleftrightarrow\quad
 g=\sum_i n_iB^{w_i},\qquad 0\le n_i<b.
\]
The tuple is unique.  Put \(z_i=n_i+1\), so \(1\le z_i\le b\).
If \(h_s\) is the coefficient of \(T^s\) in
\(\mathcal F(T,Y_z(T))\), the sum \(z=\sum_i z_i\le31b\) gives
\[
 |h_s|\le (M+1)z^2+Mz+M+1<2048Mb^2=B/4.
\]
The polynomial degree is at most \(K_c\).  Write the signed expansion
\[
 J_c=\sum_{s=0}^{K_c}(2h_s+\epsilon_s)B^s,
 \qquad \epsilon_s=\begin{cases}1&s\in\{q_1,\ldots,q_{22}\},\\0&\text{otherwise.}\end{cases}
\]
Ordinary Euclidean digit normalization has incoming carry always zero
or minus one, by induction from the strict coefficient bound.  At a target,
the normalized digit avoids the binary mask \(B-2\) exactly when
\(h_s=0\): a zero coefficient gives digit zero or one; a positive
coefficient gives a digit at least two; a negative coefficient gives a
normalized digit at least two after borrowing.  This argument explicitly
permits an incoming borrow.

Concatenation at the \(N_0\) boundary and the extraction lemma prove
\begin{equation}
 S_c\mathbin{\&}T_c=0
 \quad\Longleftrightarrow\quad
 \left\{
 \begin{array}{l}
 g=\sum_i n_iB^{w_i},\quad0\le n_i<b,\\
 r_j(n_1+1,\ldots,n_{31}+1)=0\quad(1\le j\le22).
 \end{array}\right.
 \label{eq:channel-zero}
\end{equation}
Bitwise AND is used only in the proof, never in the final equations.

For completeness, any positive parent solution gives nonnegative digits
\(n_i=z_i-1\).  They are not all zero: the fifth shifted residual equals
one at the zero vector.  Choose any sufficiently large power of two
\(b>\max_i n_i\).  All weights are positive, so the resulting
\(g\ge B\ge b\), while \(g<bB^W\), and every carry test passes.
These witnesses exist at arbitrarily large admissible powers of two;
the ten-integer construction may restrict them to \(b=4^s\).

\section{Nine positive unknowns at degree 2,628,944}\label{sec:nine}

We compose the degree-separated source of Section~\ref{sec:channels} with
an exact binary-population calibration and the inherited Jones ratio
construction. The resulting equation has nine strictly positive unknowns,
exact total degree $2\,628\,944$, and 499 live binary arithmetic operations.
Its separately translated nonnegative version has the same degree and
501 operations. The arithmetic proof below identifies the exact zero set;
no intermediate register is an additional unknown.

\subsection{Source registers and the binary threshold}
Supply the nine positive coordinates
\[
 (\epsilon_+,g,h,s,w,\phi_+,i,j_+,n_+)\in\Pos^9,
 \qquad \phi=\phi_+-1,\quad j=j_+-1,\quad n=n_+-1.
\]
Put $b=\epsilon_++1$, $a_\star=2^{1182968}$, and $B=a_\star b^2$.
Use the source registers $S_c,T_c$ from Section~\ref{sec:channels},
including its raw positive word and positive carrier.
Put $c_\star=2^{9463734}$ and define
\begin{equation}\label{eq:nine-source}
\begin{aligned}
 L&=4bB^{2241},& N&=c_\star bB^{1178},\\
 R&=(S_c+T_c)L+L-1-S_c,& G_{\rm cap}&=32bB^{108}.
\end{aligned}
\end{equation}
All exponents are fixed nonnegative integers. In particular, $N$ and the
binary block length $L$ are different polynomials.

The source theorem gives $S_c>0$ and $T_c\geq0$ on \emph{every} outer
assignment, before recovering a binary base or a bound on $g$. Hence
\[
 R=(L-1)S_c+LT_c+L-1\geq8,
 \qquad N\geq8,\qquad 0<b\leq R,N,\qquad G_{\rm cap}>0.
\]
If subsequently $8g<G_{\rm cap}$, then $g<4bB^{108}$, and the source
bounds give $S_c,T_c<B^{2241}$ and $S_c+T_c<L$.

Suppose also $b=2^k$, and write $B=2^v$, so $v=1182968+2k$.
The low mask occupies 109 base-$B$ cells, with 31 allowed digit windows;
the 22 output masks each contribute $v-1$ bits. Thus
\[
 \operatorname{pop}(T_c)=131v-31k-22,
 \qquad \log_2L=2+k+2241v.
\]
The two disjoint binary blocks of $R$ give the exact identity
\begin{align*}
 \operatorname{pop}(R)
 &=\log_2L+\operatorname{pop}(T_c)-\delta,\\
 \delta&=\operatorname{pop}(S_c)+\operatorname{pop}(T_c)
                  -\operatorname{pop}(S_c+T_c)\geq0.
\end{align*}
Here $\delta=0$ precisely when $S_c$ and $T_c$ have disjoint binary bits.
Furthermore,
\begin{align*}
 \log_2L+\operatorname{pop}(T_c)
 &=2372v-30k-20\\
 &=2\bigl(8\cdot1182968-10+k+1178v\bigr)=2\log_2N.
\end{align*}
Legendre's formula therefore proves
\begin{equation}\label{eq:nine-threshold}
 N^2\mid\binom{2R}{R}
 \quad\Longleftrightarrow\quad S_c\mathbin{\&}T_c=0.
\end{equation}
The binomial coefficient and bitwise symbol occur only in the proof.
There is no leading $L^2$ term in $R$; its omission is justified by the
unconditional positivity and exact block calculation above.

\subsection{Complete polynomial definition}
The following are polynomial registers. Letters reused from the source
construction are local to this subsection:
\begin{align*}
 U&=N^2w,& Y&=N^2s,& M_0&=RY,\\
 P&=2M_0^2U,& K&=R+1+h(P-1),& A&=M_0(U+1),\\
 B_1&=2R+1,& W_1&=bw,& C&=B_1+W_1+\phi,\\
 J&=4A-5,& Q&=3W_1C-2(W_1^2-1),\\
 D&=(A^2-1)C^2+1,& E&=iJDC^2,\\
 F&=(A^2-1)E^2+1,& G&=A+F(CD+1-A),\\
 H&=C+F(B_1+jC),& I&=(G^2-1)H^2+1.
\end{align*}
Set
\begin{align*}
 V_0&=PK,& a_0&=V_0^2-K^2+1,& a_1&=64DFI,\\
 \mu&=\bigl(K^2-4(C-KY)^2\bigr)G_{\rm cap}-8gK^2,\\
 t_0&=iJC,& T_0&=t_0(8D^2-4D-1),& V_1&=T_0GH,\\
 \Omega&=J\mu,& E_0&=Jn+Q,&
 \Theta&=E_0-\Omega(V_0+V_1).
\end{align*}
The promised equation is
\begin{equation}\label{eq:nine-final}
 \boxed{\displaystyle
 P_9=\bigl(\Theta^2-\Omega^2(a_0+a_1)\bigr)^2
                  -4\Omega^4a_0a_1=0.}
\end{equation}
The factor 64 belongs to the present definition of $a_1$; the inherited
ratio lemmas use the unscaled square predicate for $DFI$.

\subsection{The ratio interface and an exact quotient change}
The supplied \texttt{NineRatio.lean} and its accompanying mathematical
proof give the following numerical contract. Under the unconditional
preconditions established above, the four conditions
\[
 a_0\text{ is a square},\quad DFI\text{ is a square},\quad
 J\mid Q,\quad \mu>0                                      \tag{R}
\]
recover $b$ as a power of two, $8g<G_{\rm cap}$, and the binomial
condition in~\eqref{eq:nine-threshold}. Conversely, the binomial condition,
a binary $b$, and the stronger bound $32g<G_{\rm cap}$ supply the
remaining ratio witnesses. The change from the original ratio variables
to~(R) is exact, as follows.

Originally $H_{\rm old}=B_1+j_{\rm old}C$, with $j_{\rm old}>0$, and the
divisibility condition is
\[
 JF\mid QF+J(H_{\rm old}-C).
\]
Since $E$ is divisible by both $J$ and $C$, one has
$F\equiv1\pmod J$ and $F\equiv1\pmod C$. The displayed divisibility is
equivalent to $J\mid Q$ and $F\mid H_{\rm old}-C$.
Write $z=(H_{\rm old}-C)/F$. Then $z>0$ and $z\equiv B_1\pmod C$.
Because $0<B_1<C$, there is a unique $j\geq0$ with $z=B_1+jC$.
This gives exactly $H=C+F(B_1+jC)$. Conversely,
\[
 j_{\rm old}=1+\frac{F-1}{C}B_1+Fj>0
\]
is an integer and reconstructs the original $H_{\rm old}$.
The division occurs only in this proof of a witness bijection. Neither
an added variable nor an auxiliary cofinality assertion is needed.

\subsection{A polynomial majorant and recovery in the correct order}
Every register needed below is positive, with $A\geq2$, $J\geq3$, $K>1$,
$D\geq A^2$, and $H>C>1$. Also $Q>0$. Define the polynomial
\[
 \mathcal E=t_0^2(8D+1)-64D.
\]
The identities $E=t_0DC$ and $(A^2-1)C^2=D-1$ give
\begin{equation}\label{eq:nine-majorant}
\begin{aligned}
 F&=t_0^2D^2(D-1)+1,\\
 T_0^2-64DF&=\mathcal E>0,\\
 V_1^2-a_1&=\mathcal E(GH)^2+64DF(H^2-1)>0.
\end{aligned}
\end{equation}
Indeed $t_0\geq3$, so $\mathcal E\geq9(8D+1)-64D>0$.
Together with $V_0^2-a_0=K^2-1>0$, these prove
\[
 0<\sqrt{a_0}<V_0,\qquad 0<\sqrt{a_1}<V_1
\]
without assuming either radicand is a square. Moreover $J>W_1$,
$0<Q/J<3C$, and $V_1>3C$; hence $V_1>Q/J$.
For example, $U=N^2w\geq2W_1$ implies $J>W_1$, while
$T_0>t_0\geq3C$ proves the latter majorant bound.

\begin{lemma}[Exact finalizer]
For fixed outer coordinates in the stated domain, there is an $n\geq0$
satisfying~\eqref{eq:nine-final} if and only if the four conditions~(R) hold.
\end{lemma}
\begin{proof}
If $\mu=0$, then $\Omega=0$ and $P_9=(Jn+Q)^4>0$. Otherwise the real
factorization of the two-root norm implies, for some signs, that
\[
 n+\frac QJ
 =\mu\bigl(V_0+V_1\pm\sqrt{a_0}\pm\sqrt{a_1}\bigr).       \tag{*}
\]
The parenthesized expression is positive by the strict majorants; the
left side is positive because $n\geq0$ and $Q>0$. Thus $\mu>0$ is obtained
\emph{before} any integrality claim about the radicals.

The margin inequality now gives $|C-KY|<K/2$, whence $C>K$ because
$Y\geq64$. Also $P<2A^2$, $F\geq A^2$, and $I\geq1$. Consequently
\[
 a_1\geq64(A^2-1)A^2C^2
       \geq48A^4C^2>4A^4K^2>P^2K^2>a_0.
\]
Let $z=\Theta/\Omega$. The vanishing norm says
$z=\pm\sqrt{a_0}\pm\sqrt{a_1}$. This rational number is nonzero because
$a_1>a_0$. Squaring after isolating either root gives, with a suitable sign,
\[
 2z(\pm\sqrt{a_0})=z^2+a_0-a_1.
\]
Both roots are therefore rational, hence integral. In particular $DFI$
is a square, since a rational square root of an integer is integral.
Equation~(*) now gives $Q/J\in\Z$, proving~(R).

Conversely, under~(R), choose the positive integer square roots and set
\[
 n=\mu\bigl(V_0+V_1+\sqrt{a_0}+\sqrt{a_1}\bigr)-Q/J.
\]
This is an integer and is positive, since $\mu\geq1$ and $V_1>Q/J$.
One real factor of the norm vanishes, proving $P_9=0$.
\end{proof}

The ratio contract, the source theorem, and this lemma prove soundness.
For completeness, encode a positive parent solution by digits $n_i=z_i-1$.
The shifted parent has no zero at the zero vector, so its code is positive.
Choose any sufficiently large binary $b$ exceeding all digits. The source
then supplies $g<bB^{108}$, hence $32g<G_{\rm cap}$, and the binomial
condition. Apply ratio completeness, the exact quotient change, and the
lemma. Composing with the pinned compiler and K215 calibration proves
\[
 \mathrm{ZFC}\ \vdash\
 \bigl(\neg\exists\mathbf x\in\Pos^9\;P_9(\mathbf x)=0\bigr)
 \ \longleftrightarrow\ \Con.
\]
For the nine-natural version use independent coordinates
$\epsilon,g_0,h_0,s_0,w_0,\phi,i_0,j,n\in\N$, put $b=\epsilon+2$, and
replace $g,h,s,w,i$ by $g_0+1,h_0+1,s_0+1,w_0+1,i_0+1$.
These are exact domain translations; the positive polynomial is not
silently reinterpreted over unrestricted integers.

\subsection{Exact degree and evaluation cost}
The apparent top degree of the norm cancels. The supplied degree-certificate
circuit resolves that cancellation using~\eqref{eq:nine-majorant}. Write
$\Delta_1=V_1^2-a_1$. The exact polynomial identity is
\begin{align*}
 \Theta^2-\Omega^2(a_0+a_1)
 ={}&E_0^2-2E_0\Omega(V_0+V_1)\\
 &+\Omega^2\bigl(K^2-1+\Delta_1+2V_0V_1\bigr),\\
 \Delta_1={}&\mathcal E(GH)^2+64DF(H^2-1).
\end{align*}
Let $\rho=\deg R=8963$ and $\nu=\deg N=2357$. Propagation of total-degree
upper bounds through these identities gives
\[
 \deg P_9=188\rho+400\nu+8W+236=2\,628\,944,
 \qquad W=108.
\]
For clarity, $\deg\Omega=92\,184$, $\deg\Delta_1=1\,130\,104$, and
$\Omega^2\Delta_1$ uniquely supplies degree $1\,314\,472$ in the inner
expression. Its square has degree $2\,628\,944$; the other outer term has
degree $1\,791\,262$.

The upper bound is attained. Set every supplied positive coordinate equal
to a single indeterminate $x$. The leading coefficients of $R$ and $N$
are respectively $8a_\star^{4481}$ and $c_\star a_\star^{1178}$.
The leading coefficient of the final polynomial is exactly
\[
 \boxed{2^{742}a_\star^{1314060}c_\star^{400}>0.}
\]
This proves exact total degree, rather than a syntactic bound or a random
numerical estimate. The same coefficient survives the natural-coordinate
translations. The primary circuit has 170 additions, 25 subtractions,
and 304 multiplications, totaling 499 operations; the natural version has
175 additions, 22 subtractions, and the same 304 multiplications, totaling
501. Fixed-power evaluation is expanded into these binary gates. The
larger cancellation-resolved circuit is a degree certificate and is not
substituted for the primary evaluator's cost.

\section{Ten unrestricted integers}
\label{sec:ten}

Supply exactly ten unrestricted integer coordinates
\[
 (f,g,h,k,l,w,x,y,z,m)\in\Z^{10}.
\]
Use the source of Section~\ref{sec:channels}, with
\begin{gather*}
 b=1+3f,\qquad B=a_Sb^2,\qquad \mathcal C=bB^{108},\qquad L_c=B^{2241},\\
 X=L_c^2+(S_c+T_c)L_c+L_c-1-S_c,\\
 Y=a_S^{15}\,2^{\kappa-21}B^{2356}b.
\end{gather*}
The monomial cap is \(\mathcal C=a_S^{108}b^{217}\), as required
by the inherited Sun bridge.  The source output \(Y\) is a fixed
polynomial, not an exponential with a variable exponent.

\subsection{The arithmetic source contract}

Before recovering a binary base, the reverse bridge has \(f\ne0\),
\(b\) an integer square, and \(0\le g<2\mathcal C\).  Since
\(b\equiv1\pmod3\), these hypotheses imply \(b\ge4\) and \(f>0\).
The preceding source proof applies and gives
\(S_c,T_c\ge0\) and \(S_c+T_c<L_c\), including the stronger summed
bound~\eqref{eq:source-sum-range}.  Hence
\[
 X\ge L_c^2>3b,\qquad Y>\max(b,256).
\]
The exponent \(\kappa\) is even, so \(B\equiv N_0\equiv L_c\equiv1\pmod3\).
We have \(M_c\equiv0\), \(A_T\equiv22\), and
\(T_c\equiv-22\pmod3\).  Consequently
\[
 X\equiv1+T_c\equiv0\pmod3.
\]
These range and congruence facts do not assume \(b\) is a power of two.

After binary-base recovery, write \(b=2^s\), \(B=2^v\).  For
nonnegative \(S,T\) with \(S+T<L_c\), the three disjoint blocks in
\(L_c^2+(S+T)L_c+L_c-1-S\) give
\[
 \pop(X)=1+\log_2L_c+\pop(T_c)
       -\bigl(\pop(S_c)+\pop(T_c)-\pop(S_c+T_c)\bigr).
\]
The subtracted carry deficit is nonnegative, and vanishes exactly when
\(S_c\mathbin{\&}T_c=0\).  The exact mask population is
\(\pop(T_c)=131v-31s-22\), so
\[
 1+\log_2L_c+\pop(T_c)=2372v-31s-21=\log_2Y.
\]
Legendre's formula therefore proves
\[
 Y\mid\binom{2X}{X}\quad\Longleftrightarrow\quad
 S_c\mathbin{\&}T_c=0.
\]
Conversely a parent zero can be encoded with arbitrarily large
\(b=4^s\), putting \(f=(b-1)/3\).  It has \(b\le g<\mathcal C\)
and the required divisibility.  This is the full source contract; the
Sun bridge does not require \(Y\) to be a square or equal to \(L_c^2\).

\subsection{Literal final polynomial}

The following finalizer is inherited from the inspected package-31 and
wave-7 construction; only its source polynomials have changed.  Define
\begin{gather*}
 L=lY,\quad U=2LX,\quad V=4gwY,\quad W_b=bw,\\
 A=U(V+1),\quad B_*=2X+1,\quad C=B_*+(A-2)h,\\
 P=U^2V,\quad K=X+1+k(P-2),\quad S_*=2A-5,\\
 T_*=3W_bC-2(W_b^2-1),\quad R_2=(P^2-4)K^2+4,\\
 \Delta_*=8\mathcal C^3gK^2-g^2\{32(C-KL)^2\mathcal C^3+g^2K^2\},\\
 O=f^2l^2\Delta_*,\quad\xi=zx,\quad D=(A^2-4)C^2+4,\\
 E=C^2D\xi,\quad F=4(A^2-4)E^2+1,\\
 G_*=2(A+2)(A-2)^2E^2-1,\quad\Omega_*=16E,\\
 H_*=C+(B_*-C)F+\Omega_*Fy,\quad I_*=(G_*^2-1)H_*^2+1.
\end{gather*}
Put
\begin{align*}
 r&=T_*+S_*y,&\Sigma&=\xi S_*^5,\\
 \Theta_*&=(r^3-rS_*^2-S_*^3)(r^2+23S_*^2),\\
 a_0&=b\{1+(3O-4)b^2z^2\},&a_1&=R_2,&a_2&=DFI_*,\\
 w_0&=1,&w_1&=a_0,&w_2&=a_0(P^2+K^2+3),\\
 Z_*&=\Sigma m-\Theta_*,&A_i&=\Sigma^2w_i^2a_i\quad(0\le i\le2),\\
 V_*&=Z_*^2+A_2-A_0-A_1,&
 W_*&=V_*^2+4Z_*^2A_2-4A_0A_1.
\end{align*}
The single equation is
\begin{equation}
 \boxed{F_{10}=W_*^2-16V_*^2Z_*^2A_2=0.}
 \label{eq:ten-final}
\end{equation}
All powers in these formulas are fixed nonnegative integers.  In
particular, no square root, factorial, binomial coefficient, bit operation,
or divisibility predicate occurs in the actual polynomial.

\subsection{Why the inherited finalizer still applies}

Sun's arithmetic bridge, specialized as in the preserved proof, supplies
both directions for any source satisfying the contract above.  The reverse
direction requires \(b\) square, \(O>0\), \(\xi\ne0\),
\(R_2\) and \(DFI_*\) square, and \(S_*\mid T_*\).  It recovers
\(b\) as a power of two, \(0<g<2\mathcal C\), and the central-binomial
divisibility.  The forward direction supplies these conditions, together
with the residue interpolation needed for the finalizer.  These are the
same arithmetic bridge hypotheses as in the inspected construction.

For clarity, the finalizer's main logical steps are recorded here.
The polynomial \((v^3-v-1)(v^2+23)\) has a root modulo every positive
integer and no rational root.  Prime valuations in its homogeneous
factors show that \(\Sigma\mid\Theta_*\) implies \(S_*\mid T_*\).
The parity of \(S_*\) makes it nonzero, and the absence of rational
roots makes \(\Theta_*\ne0\).  If \(\Sigma=0\), then
\(F_{10}=\Theta_*^8\ne0\), excluding that branch.

For \(z\ne0\), the factors \(b\) and
\(1+(3O-4)b^2z^2\) are coprime.  If their product is a square,
they are signed squares; \(b\equiv1\pmod3\) excludes a negative
square and the zero-factor exception.  The positivity Pell lemma then
gives \(O>0\).  Conversely, the forward bridge has \(O\ge4\);
a suitable positive Pell unit makes \(a_0\) a square.  Residue
interpolation supplies a multiple \(\xi=zx\) and the required
class of \(y\), making \(a_2\) a square and
\(\Sigma\mid\Theta_*\).

The structural weights satisfy
\[
 |w_1|>\sqrt{|a_0|},\qquad
 |w_2|>\sqrt{|a_0|}+|w_1|\sqrt{|a_1|}.
\]
Here \(|a_0|>1\) whenever \(z\ne0\), and
\(\sqrt{|a_1|}\le P^2+K^2+2\).  In the multiquadratic field,
any nontrivial sign-changing automorphism of a weighted root sum has
a largest changed index whose magnitude exceeds the sum of all smaller
changed magnitudes.  Thus a rational weighted sum forces every radicand
to be an integer square.  The exact identity
\[
 F_{10}=\prod_{\epsilon\in\{\pm1\}^3}
 \left(Z_*+\sum_{i=0}^2\epsilon_i\sqrt{A_i}\right)
\]
now proves the finalizer equivalence.  This argument also handles
negative radicands using complex absolute values, and zero radicands
separately.  At a zero, the bridge and~\eqref{eq:channel-zero} recover
a parent solution.  From a parent solution, choose the forward bridge
witnesses and set \(m\) to cancel one integral signed factor.

The full inherited interpolation, local-root, and bridge proof is kept
in \code{upstream/wave7_sun_popcount.tex}, with its precise primary
references~\cite{sun2021,sun2026}.  The root verifier checks the norm
identity by exact polynomial multiplication; the unbounded bridge is
an identified mathematical dependency.

\subsection{Exact degree and circuit cost}

The source registers have degrees
\[
 \deg X=8964,\qquad\deg Y=4713,\qquad\deg\mathcal C=217.
\]
Representative downstream degrees are
\[
\begin{array}{c|rrrrrrr}
 &A&C&P&K&D&E&F\\\hline
\deg&18393&18394&32071&32072&73574&110364&257514
\end{array}
\]
and
\[
\begin{array}{c|rrrrr}
 &G_*&H_*&I_*&\Sigma&\Theta_*\\\hline
\deg&275907&367879&1287572&91967&91980.
\end{array}
\]
The exact leading-coefficient certificate in
\code{generated/ten_integer_degree.json} propagates the degree bound
through every gate and exhibits a strictly positive final diagonal
leader.  The independent modular leading residues in
Section~\ref{sec:status} verify it with the actual fixed coefficients.
The resulting exact degree is \(8\,317\,400\).  There are 544 live
binary operations: 174 additions, 26 subtractions, and 344 multiplications.
The baseline comparison concerns unrestricted integer zeros throughout.

\section{A quartic in forty-five unrestricted integers}
\label{sec:quartic45}

The independent low-degree contribution reduces the inspected
fifty-integer quartic to \textbf{forty-five unrestricted integer unknowns},
while retaining degree four.  It uses 28 quadratic residuals instead of
29.  Its expanded representation has 407 symbolic terms, compared with
362 in the chosen parent; the improvement concerns variable count at
fixed degree, not monomial count or evaluation cost.

This construction keeps the complete package-34 coefficient tuple
\begin{center}\code{mixed-parity-unrounded-1064921},\end{center}
with
\(B=2^{1064921}\), \(\beta=B-1\), and ports
\(\beta,K_0,M_C,M_F\).  The last port is the already shifted
\(M_{F,\mathrm{native}}+B-1\).  It is a separate complete binding from
the tuple used for the principal twelve- and thirteen-variable equations.
The four raw payloads, their finite construction, the literal K215
transition table, and the relevant inherited proofs are included in
\texttt{quartic/upstream/}.  The canonical accepting witnesses for this
tuple have \(q=2^t\) with \(t\) odd.

\subsection{The two guard reductions}

Write
\[
 \operatorname{Pos}(u,v,t)=1+u^2+v^2+t(t+1).
\]
This expression is positive on all integer triples and represents every
positive integer.  For the converse, apply the three-square theorem to
\(4n-3\): its residue modulo four forces one odd square and two even
squares, yielding the displayed representation after division by four.

\begin{lemma}[A common-sign guard]
Suppose \(q>0\) and \(C=q-F-\alpha\).  Then
\[
 F\alpha>0\quad\hbox{and}\quad FC>0
 \quad\Longleftrightarrow\quad F,\alpha,C>0.
\]
\end{lemma}
\begin{proof}
The two products make the three nonzero integers have the same sign.
Their sum is \(q>0\), so that sign is positive.  The converse is immediate.
\end{proof}

\begin{lemma}[A dyadic guard]
For \(\beta\ge3\), the equation
\(\beta J+1=2v^2\), with \(J,v\in\mathbb Z\), implies \(J>0\).
Every canonical radix \(q=2^t\) with \(t\) odd satisfies this equation
with \(v=2^{(t-1)/2}\).
\end{lemma}
\begin{proof}
If \(J\le-1\), its left side is negative; if \(J=0\), it is one,
which cannot be twice an integer square.  The second assertion is a
direct substitution.
\end{proof}

\subsection{The complete quadratic system}

The 32 arithmetic coordinates are
\[
\begin{gathered}
J,F,\alpha,z,f,h,r,v_s,w,\tau,\eta,\zeta,y,\gamma,o,j,\\
N,E,P,c,a,X,D,C_2,L,F_2,V,g,I,B_{\rm aux},\Delta,u.
\end{gathered}
\]
Add one coordinate \(v_q\) and four triples
\(\boldsymbol e,\boldsymbol\xi,\boldsymbol b,\boldsymbol d\).
All 45 coordinates range over \(\mathbb Z\).  In the following formulas,
the letters
\[
\begin{gathered}
 q=\beta J+1,\quad C=q-F-\alpha,\quad
 k=\eta+\zeta,\quad H=4a+3,\\
 R=(q-F)(N-q)-CL+(M_C+qM_F)J
\end{gathered}
\]
are polynomial abbreviations, not additional coordinates.  Retain the
following 23 arithmetic equations:
\begin{alignat*}{2}
 L&=q^2-1,&\qquad N&=q(L+1),\\
 X&=wq,&a&=E+(2v_s+1)N,\\
 E&=X(a-E),&c&=k(a-E)+\eta,\\
 P&=E(c-\eta)+k,&\tau^2&=P(P-k)+1,\\
 D&=X+ac+\gamma H,&\Delta&=a^2+H,\\
 C_2&=c^2,&D^2&=\Delta C_2+1,\\
 k&=R+1+hE,&(K_0+w)C+q-F-z(q-1)&=1,\\
 F_2&=f^2,&V&=of-c,\\
 V-cj+k-1-hE&=0,&g&=y^2-V^2,\\
 I&=(2r+1)C_2,&B_{\rm aux}&=I^2,\\
 B_{\rm aux}&=\Delta(F_2-1),&y^2-B_{\rm aux}g&=1,\\
 w&=u^2.&&
\end{alignat*}
Add the five quadratic guards
\begin{align*}
 \eta&=\operatorname{Pos}(\boldsymbol e),&
 \zeta&=\operatorname{Pos}(\boldsymbol\xi),&
 q&=2v_q^2,\\
 F\alpha&=\operatorname{Pos}(\boldsymbol b),&
 FC&=\operatorname{Pos}(\boldsymbol d).
\end{align*}
Here \(\operatorname{Pos}(\boldsymbol e)\) means evaluation at its three
coordinates.  Every displayed residual is quadratic.  If they are
\(r_1,\ldots,r_{28}\), define the single equation
\begin{equation}
 \boxed{\displaystyle Q_{45}=\sum_{i=1}^{28}r_i^2=0.}
 \label{eq:quartic45}
\end{equation}
The coefficient of \(\tau^4\) is one, so the degree is exactly four.
The literal ordered residuals, variable list, and exact expansion are
\texttt{quartic/quartic45.json} and
\texttt{quartic/quartic45\_expanded.json}.

\subsection{Why the unguarded transport quotient is harmless}

The old fifty-variable system had independent ternary guards for
\(J,F,\alpha,z,\eta,\zeta\).  Delete the first four guards and their
twelve private variables; retain the last two; add \(v_q\) and the two
new triples.  Thus the variable count is
\(50-12+7=45\), and the residual count is \(29-4+3=28\).
Because these are domain changes, an algebraic identity alone would not
prove the required existential equivalence.

At a zero of \eqref{eq:quartic45}, the dyadic guard gives \(J>0\),
and the common-sign guard gives \(F,\alpha,C>0\).  The retained
guards give \(\eta,\zeta>0\).  The only missing original positivity
condition is \(z>0\).  Its recovery must respect the proof order.

In the inherited signed Pell proof, the early use of \(z>0\) was
solely to infer \(C\ge0\) from transport.  We already have the stronger
conclusion \(C>0\), independently.  Hence the same packing bounds follow:
\[
 q(q^2-1)<R<q^4-2,\qquad R-2>3q+1.
\]
The square-scale exclusion gives \(w>0\).  The existing odd-scale
argument, with \(s=2v_s+1\), excludes negative \(s\); it uses the
Pell equations, these bounds, and the exposed congruences
\(V\equiv-c\pmod{|f|}\), \(V\equiv-R\pmod{|c|}\), without
another use of \(z\).  Positive-scale rank recovery and the signed
main-root argument then give
\[
 X=2^R.
\]
For a detailed dependency trace see the included
\texttt{ODD\_SCALE\_27\_PROOF.md}, Sections 1--4, and
\texttt{SIGNED\_DOMAIN\_29.md}, through its signed-gamma and
ordinary-strong arguments.  This explicitly supplies \(C>0\) in place
of the original preliminary transport inference; it does not invoke
the final signed-domain theorem before restoring its guard.

Since \(R>3q\), we have \(w=2^R/q>q>F\).  The retained transport
equation now reads
\[
 (K_0+w)C=(z-1)(q-1)+F.
\]
Its left side exceeds \(F\), so \(z>1\).  All original positivity
conditions are therefore restored.  Choosing their ternary representations
produces a zero of the old fifty-variable quartic, and its inherited
compiler theorem yields a finite ZFC refutation.

Conversely, canonical accepting witnesses have \(q=2^t\) with odd
\(t\), and \(F,\alpha,C>0\).  Set \(v_q=2^{(t-1)/2}\) and use
the positive representations of \(F\alpha\) and \(FC\).  Every new
equation holds.  Consequently the ordinary mathematical argument gives,
formalizably in ZFC,
\[
 \text{there is no integer zero of }Q_{45}
 \quad\Longleftrightarrow\quad \operatorname{Con}(\mathrm{ZFC}).
\]

The portable finite audit checks all 28 residual degrees, the 45 live
coordinates, the exact expansion, \([\tau^4]Q_{45}=1\), the complete
fixed payloads, and the pinned inherited bytes.  The unbounded Pell and
compiler arguments remain explicit mathematical dependencies; this
section does not claim a new end-to-end Lean/Rocq certificate or an
exported object-language ZFC derivation.


\section{Verification, reproduction, and formalization status}
\label{sec:status}

\subsection{What the finite audits establish}

The new proofs separate unbounded implications from finite certificates.
The source checker reconstructs all 22 target identities as polynomial
identities with four indeterminate coefficient ports, and exhausts all
23,599 relevant source/expansion interactions.  The toy source checker
deliberately allows a linear and a quadratic term to share a reversed
position.  It checks 1,020 source/carry cases, including three genuine
zero cases, 680 forbidden-digit controls, and 36 unconditional-positivity
regressions.  These diagnostics would be meaningless as a proof of ZFC
consistency; their purpose is to expose implementation mistakes in the
source transformation.

The root verifier reads the emitted circuits without importing their
builder.  It authenticates all native and derived coefficient bytes,
checks every gate reference, all variable leaves, liveness, and operation
counts.  It compares the complete output and key named registers with
independent literal formulas at 72 deterministic modular assignments per
equation.  Four formal polynomial identities check the Pell majorant,
its defect, the nine-variable degree cancellation, and the eight-factor
three-radical norm.  The latter norm has 35 terms.

The exact-degree check propagates a multivariate upper bound and the
coefficient at that degree on a diagonal.  The cancellation-resolved
identity is used for nine variables.  With the actual coefficient
payloads, the independently obtained leading residues are
\[
\begin{array}{c|rrr}
\text{equation}&101&1009&1000003\\\hline
F_9&77&467&509827\\
F_{10}&80&969&99517
\end{array}
\]
Any nonzero residue proves the required lower bound.  Thus the reported
degrees are exact, not syntactic upper bounds or estimates from sampled
function values.  The low-degree checker evaluates full univariate
polynomials with the actual tuple, using \(t=3x\) and all other
coordinates equal to \(x\) to avoid a vanishing top form.  The universal
upper bounds are supplied by the cancellation identities in its proof.
The quartic checker expands the sum of squares exactly and verifies its
degree-four coefficient, all 45 live coordinates, 28 residuals, and
the separate coefficient binding.

\subsection{Reproduction commands}

All computational audits use the Python standard library and can be run
from the packet root:
\begin{verbatim}
python3 build.py
python3 materialize.py
python3 verify.py
python3 data/verify_channels.py
python3 data/check_toy_channels.py
python3 coefficients/audit_coefficients.py
python3 low_degree/verify.py
python3 quartic/build_verify.py
\end{verbatim}
The all-in-one runner \code{reproduce.py} also checks the pinned local
sources and assembles a machine-readable receipt.  The low-degree builders
and both coefficient builders are supplied for regeneration.  The source
position search is not required: the small fixed certificate is the proof
input.  The two source checkers explicitly reject Python optimization
flags, because their internal assertions are part of their audit.

Compile the paper from \code{report/} with two passes of
\code{pdflatex -halt-on-error main.tex}.  The final packet includes the
compiled PDF and all its TeX inputs.  \code{MANIFEST.sha256} records the
distributed files.  An extracted copy was used for the final reproduction
check, so the executable audits do not depend on the development checkout.

\subsection{Ordinary provability versus an exported object-language proof}

The intended consistency predicate is syntactic: there is no finite
derivation of falsity from the recursively presented axioms of ZFC.
The accepted K215 proof states that its blank halting is equivalent to
the existence of such a derivation, by ordinary reasoning formalizable
in ZFC.  The present work composes that inherited theorem with explicit
new polynomial reductions.  Integer arithmetic, finite coefficient
extraction, Pell recurrences, exact divisibility, and the witness
constructions used here can all be developed internally in ZFC; none
requires assuming \(\Con\).

This mathematical claim must be distinguished from the repository's
remaining proof-assistant obligation.  At the inspected revision,
\code{Computability/DiophantineConsistency/OBJECT_ZFC.md} states that
an external Lean theorem about \code{BProv ZFCax_s [] Form.fBot} is
not itself an object-language derivation of the equivalence.  The concrete
target sentences and a complete derivation of their equivalence are not
yet assembled.  A formal proof must handle the model's own arithmetic,
including nonstandard witnesses, and cannot use numeral agreement or
bounded consistency as a substitute.  We preserve this boundary: the
archive supplies new mathematical proofs and exact executable finite
certificates, not a new kernel-checked end-to-end theorem or a completed
\code{BProv} endpoint.

Neither the existence nor the absence of solutions is asserted
unconditionally in this paper.  The equivalences are the result.

\appendix
\section{Complete channel polynomials}
\label{app:channels}

The following table specifies every term of the three channel polynomials
used in Section~\ref{sec:channels}.  A row \((d,e,c)\) contributes
\(cT^e\) to \(A_d(T)\).  The four symbols in its coefficient column
are the fixed native tuple of Section~\ref{sec:coefficients}.
Thus the table, the fixed payloads, the weights and the displayed register
formulas give an explicit polynomial without an appeal to a general
existence theorem for Diophantine representations.
\small
\begin{longtable}{@{}rrl@{}}
\toprule Degree channel $d$ & Exponent $e$ & Coefficient $c$\\\midrule
\endfirsthead
\toprule Degree channel $d$ & Exponent $e$ & Coefficient $c$\\\midrule
\endhead
\bottomrule\endfoot
0 & 39 & $-2$\\
0 & 238 & $2K_0$\\
0 & 368 & $-2$\\
0 & 1093 & $-6$\\
0 & 1147 & $-2$\\
0 & 1388 & $-2$\\
0 & 1621 & $2$\\
0 & 1861 & $-2$\\
1 & 0 & $2$\\
1 & 29 & $2$\\
1 & 32 & $2$\\
1 & 134 & $2\beta+2\beta K_0$\\
1 & 140 & $2$\\
1 & 159 & $-2K_0$\\
1 & 162 & $-2-2K_0$\\
1 & 238 & $2$\\
1 & 461 & $2$\\
1 & 499 & $2$\\
1 & 539 & $-2$\\
1 & 545 & $2$\\
1 & 652 & $2$\\
1 & 732 & $2$\\
1 & 741 & $-2$\\
1 & 760 & $2$\\
1 & 828 & $2$\\
1 & 966 & $2$\\
1 & 977 & $-4\beta$\\
1 & 984 & $2$\\
1 & 1001 & $2$\\
1 & 1090 & $-8$\\
1 & 1251 & $2$\\
1 & 1268 & $-2$\\
1 & 1282 & $2$\\
1 & 1284 & $-2\beta$\\
1 & 1291 & $-2$\\
1 & 1298 & $2$\\
1 & 1362 & $2$\\
1 & 1397 & $2$\\
1 & 1404 & $-2$\\
1 & 1407 & $-2$\\
1 & 1515 & $-2$\\
1 & 1517 & $-2M_F-2M_C+4\beta$\\
1 & 1524 & $2$\\
1 & 1545 & $-2$\\
1 & 1582 & $-2$\\
1 & 1671 & $2$\\
1 & 1877 & $2$\\
1 & 1882 & $-6$\\
1 & 1910 & $-2$\\
2 & 2 & $-1$\\
2 & 9 & $-1$\\
2 & 36 & $\beta$\\
2 & 61 & $-1$\\
2 & 64 & $-1$\\
2 & 74 & $-\beta$\\
2 & 111 & $-2$\\
2 & 280 & $2$\\
2 & 283 & $-1$\\
2 & 337 & $-1$\\
2 & 514 & $1$\\
2 & 517 & $1$\\
2 & 533 & $-1$\\
2 & 536 & $-1$\\
2 & 543 & $-2$\\
2 & 603 & $-1$\\
2 & 819 & $-2$\\
2 & 824 & $1$\\
2 & 843 & $-1$\\
2 & 864 & $-1$\\
2 & 873 & $-2\beta^{2}$\\
2 & 1087 & $-2$\\
2 & 1113 & $-2$\\
2 & 1120 & $1$\\
2 & 1123 & $1$\\
2 & 1125 & $2$\\
2 & 1144 & $-1$\\
2 & 1164 & $-\beta$\\
2 & 1187 & $-\beta$\\
2 & 1385 & $1$\\
2 & 1391 & $-1$\\
2 & 1411 & $-\beta$\\
2 & 1413 & $-2\beta M_F+2\beta^{2}$\\
2 & 1420 & $\beta$\\
2 & 1439 & $1$\\
2 & 1441 & $-\beta$\\
2 & 1445 & $-1$\\
2 & 1448 & $-1$\\
2 & 1584 & $1$\\
2 & 1623 & $-2$\\
2 & 1633 & $2$\\
2 & 1680 & $-1$\\
2 & 1787 & $2$\\
2 & 1879 & $-4$\\
2 & 1910 & $-1$\\
\end{longtable}
\normalsize


\begin{thebibliography}{9}
\raggedright
\bibitem{repo}
V. Reshetnikov and contributors, \emph{ProveIt}, inspected revision
\code{4b4349064a750ba6347293b55874a7f92abdd0db}, especially
\code{Computability/DiophantineConsistency/con-zfc-diophantine-forms}
and \code{docs/incoming}.
\url{https://github.com/VladimirReshetnikov/ProveIt/tree/4b4349064a750ba6347293b55874a7f92abdd0db}.

\bibitem{jones}
J. P. Jones, \emph{Universal Diophantine equation}, Journal of Symbolic
Logic \textbf{47} (1982), 549--571.  The numerical ratio interface used
here is preserved in \code{upstream/NineRatio.lean}, with its supporting
\code{Ratio.lean} and \code{ProductPell.lean} source.

\bibitem{sun2021}
Z.-W. Sun, \emph{Further results on Hilbert's Tenth Problem},
Science China Mathematics \textbf{64} (2021), 281--306.
\url{https://arxiv.org/html/1704.03504v7}.

\bibitem{sun2026}
Z.-W. Sun, \emph{Ten unknowns for Hilbert's tenth problem over the integers},
2026 preprint, version 2.  The inherited ten-integer bridge is specialized
from Propositions 6.1--6.2.
\url{https://arxiv.org/html/2609.01594v2}.

\bibitem{formal}
J. Bayer, M. David, M. Hassler, Yu. Matiyasevich, and D. Schleicher,
\emph{Diophantine equations over \(\Z\): universal bounds and parallel
formalization}, 2025 preprint.
\url{https://arxiv.org/html/2506.20909v2}.
\end{thebibliography}
\end{document}
